El distribuir \emph{software} de las cocinas de los ingenieros a las mesas de los usuarios finales no ha sido una tarea sencilla~\cite{mantylaSoftwareDeploymentActivities2011}. Los artefactos de \emph{software} se comportan igual a las plantas exóticas cuando son trasplantadas a un hábitat diferente al que están acostumbradas: se marchitan~\cite{Dolstra2006}. Como las plantas, el \emph{software} ``crece y evoluciona'' en el \emph{hardware}, sistema operativo y librerías específicas de la computadora del ingeniero, pero en el momento que esos artefactos son llevados a los ecosistemas extraños, que son los dispositivos de los usuarios finales, el que funcione o no se vuelve una apuesta ante la que no se tiene control\dots{} ¿o sí?~\cite{Dolstra2006}.

El problema anteriormente descrito es el sujeto de estudio del campo de ``Manejo de Configuración de \emph{Software}'' (o SCM por sus siglas en inglés), donde se reconoce que la ejecución correcta de \emph{software} no solamente depende de su código fuente, sino del contexto que lo rodea~\cite{Dolstra2006}. Con los años se ha desarrollado una familia de \emph{software}, llamada \textbf{manejadores de paquetes}, responsable de modificar el entorno global de las computadoras objetivo para conseguir las condiciones ideales para cada aplicación. Al día de hoy se han convertido en una familia tan variada que se han vuelto una característica diferenciadora en las diferentes distribuciones de Linux o lenguajes de programación~\cite{Gibb2026}. Una corriente opuesta es la \textbf{virtualización}, que consiste en empaquetar las aplicaciones junto a los entornos completos que necesitan y ejecutarlas de forma aislada; soluciones de este tipo son muy usadas en servicios de la nube~\cite{PDFInfrastructureCode}.

Sin embargo, como se ilustra en la Figura~\ref{figura1}, los manejadores de paquetes lidian con el problema de que intentar satisfacer a varias aplicaciones en un entorno global puede llevar a conflictos irresolubles: cuando \texttt{FOO} y \texttt{BAR} dependen de versiones distintas de Node, el entorno global obliga a elegir una sola versión compatible (Node v23.8). La virtualización elimina ese conflicto permitiendo que cada aplicación lleve su propia versión en entornos separados, pero a cambio puede duplicar dependencias compartidas como Clang 19.2~\cite{Zwinger2026}, incrementando el consumo de almacenamiento en varios casos~\cite{zhaoEndtoendHighperformanceDeduplication2024a}.

\begin{figure}[htbp]
\centering
\includegraphics[width=\textwidth]{Figure1}
\caption{Dependencias de dos programas ficticios \texttt{FOO} y \texttt{BAR} en manejadores de paquetes vs. virtualización.}
\label{figura1}
\end{figure}

\emph{La unión hace la fuerza}, dando origen en 2003 a Nix como una tercera alternativa que fusiona ideas de ambas corrientes partiendo de la idea de que usar los mismos ingredientes y pasos debería producir el mismo resultado sin importar la computadora. Nix garantiza lo primero mediante identificadores (ID) únicos en un entorno aislado (\emph{Nix Store}) que, como se ve en la Figura~\ref{figura2}, permiten la coexistencia de versiones distintas del mismo paquete y la reutilización de dependencias compartidas, resolviendo los problemas previos de \texttt{FOO} y \texttt{BAR}; y lo segundo mediante entornos aislados que aseguran la reproducibilidad~\cite{Dolstra2006}. La elegancia de Nix reside en cómo construye estos identificadores y entornos aislados.

\begin{figure}[htbp]
\centering
\includegraphics[width=0.7\textwidth]{Figure2}
\caption{Manejo de paquetes en Nix con identificadores únicos permite coexistencia y reutilización de paquetes.}
\label{figura2}
\end{figure}

Fue este enfoque centrado en seguir recetas explícitas el que permitió a Nix conseguir una serie de hitos importantes al contar con uno de los repositorios de paquetes generales más grandes de Linux~\cite{marakasovRepositoryStatistics}, de los cuales 700 mil han demostrado poder replicarse de forma binariamente idéntica en diferentes computadoras~\cite{Malka2025}. Aconteció que muchas de las ideas podían generalizarse hasta el punto de reproducir casi por completo un sistema operativo, dando origen a la distribución NixOS~\cite{Dolstra2008}.

A pesar de ello, Nix ha gozado de una adopción bastante reducida en comparación a las otras herramientas discutidas~\cite{stackoverflowMostPopularTechnologies} ¿Cuál es entonces su talón de Aquiles? Tal parece que no son necesariamente problemas técnicos, sino de experiencia de uso; en encuestas hechas en el foro oficial, la comunidad resalta problemas importantes con la documentación, errores crípticos y un \emph{Domain Specific Language} (DSL) difícil de dominar~\cite{2022NixSurvey2022,2025CommunitySurvey,NixCommunitySurvey2023}; además, en otra encuesta, se estima que los usuarios perciben requerir un tiempo de 5 años para dominar la herramienta a pesar de que la mayoría lo usa a diario. Llevando a un raro caso donde a pesar de que a la comunidad le encanta la idea detrás de Nix~\cite{NixCommunitySurvey2024}, sus problemas de usabilidad son tan severos que podrían estar impidiendo su uso, lo que concuerda con observaciones de otros estudios en herramientas con situaciones similares~\cite{goodwinFunctionalityUsability1987}.

El concepto de que los desarrolladores también son usuarios dio origen al campo de estudio de Experiencia de Desarrollo (o DX por sus siglas en inglés), donde el estudio sobre cómo los desarrolladores perciben sus herramientas ha sido un tema frecuente~\cite{Razzaq2024} sobre el que ya se han desarrollado algunos instrumentos como DEXI para evaluar dichas dimensiones~\cite{Kuusinen2016}. Y dado el trayecto de intentos por mejorar la DX en Nix~\cite{caddetNixNickel,gagarinFourMonthsNix,hufschmittCurrentStatePtyx,fricklerhandwerk2022}, el presente trabajo busca ser una aplicación de las técnicas aprendidas en el campo de DX, en conjunto con el diseño de lenguajes, para evaluar si un Lenguaje de Propósito Específico Embebido (eDSL por sus siglas en inglés) en TypeScript~\cite{Typescript} podría ayudar a reducir la barrera de entrada para nuevos desarrolladores en la herramienta.
